
%   Настройка основного вида документа
%	Сделано на основании Положения о ВКР
%   Документ не редактировать!
%
%   Все изменеия вносить в custom_preampule.tex
%
%   Автор: Гаврилов А. Г.


%===================================================
% Настройка шрифтов и языка
%===================================================

\usepackage{fontspec}
\setmainfont{Times New Roman}
\setsansfont{Arial}
\setmonofont{Cascadia Code SemiLight}%[Scale=0.865] 
             

\usepackage{polyglossia}
\setdefaultlanguage[babelshorthands=true]{russian}
\setotherlanguage{english}


\setlength{\parindent}{1.25cm}

%===================================================
% Важные пакеты
%===================================================


\usepackage{etoolbox} %много вспомогательных команд
\usepackage{calc} % Добавляем пакет для команды \widthof еще он добавляет возможность вычисления длин типа 1in+6cm-10pt
\usepackage{totcount} %счетчики, сохраняющие значене между компиляциями
\usepackage{varwidth} %MinipageАвтоматическойШирины \begin{varwidth} \end{varwidth}
\usepackage{expl3}
\usepackage{amsmath, amssymb}

%===================================================
% Настройка полей
%===================================================
\usepackage{geometry}
\geometry{
	a4paper,
	left=30mm,
	right=15mm,
	top=20mm,
	bottom=20mm,
}

%===================================================
% Картинки
%===================================================

\usepackage{graphicx}
\graphicspath{{images}}
\usepackage{color}
\usepackage{tikz}

%===================================================
% Таблицы
%===================================================

\usepackage{longtable} %многостраничные таблицы
\usepackage{tabularx} %таблицы с полями автоширины 
\usepackage{array}


% Шрифт в таблицах (требования к оформлению ВКР 2026: 12 пт, одинарный интервал)
\AtBeginEnvironment{table}{%
	\def\insidetable{true}%
	\renewcommand{\arraystretch}{1.0}%
}
\AtEndEnvironment{table}{\renewcommand{\arraystretch}{1.5}}

\AtBeginEnvironment{tabular}{\ifdefined\insidetable\ifDebug\color{red}\fi\fontsize{12pt}{12pt}\selectfont\fi}
\AtBeginEnvironment{longtable}{\ifdefined\insidetable\ifDebug\color{red}\fi\fontsize{12pt}{12pt}\selectfont\fi}
\AtBeginEnvironment{tabularx}{\ifdefined\insidetable\ifDebug\color{red}\fi\fontsize{12pt}{12pt}\selectfont\fi}

\renewcommand{\arraystretch}{1.5} % межстрочный интервал вне таблиц float

%===================================================
% Математические формулы
%===================================================

\usepackage{amsmath}
\usepackage{amsfonts}

% Нумерация формул в пределах раздела: (2.1) справа (требования к оформлению ВКР 2026)
\numberwithin{equation}{section}
\renewcommand{\theequation}{\thesection.\arabic{equation}}
\makeatletter
\def\tagform@#1{\maketag@@@{(\ignorespaces#1\unskip)\hfil}}
\makeatother




%===================================================
% Verbatim
%===================================================

\usepackage{verbatim}
%===================================================
% Межстрочный интерва
%===================================================

\usepackage{setspace}
%\onehalfspacing % Межстрочный интервал 1.5
\setstretch{1.425}

%===================================================
% Отключение переносов слов
%===================================================

\usepackage[none]{hyphenat}


%===================================================
% Заголовки глав и параграфов
%===================================================

\setcounter{secnumdepth}{4} %Глубина вложенности заголовокв в документе
\setcounter{tocdepth}{3} % в оглавлении: section, subsection, subsubsection (без paragraph/subparagraph)

\usepackage{titlesec} % Для настройки заголовков
\titleclass{\paragraph}{straight}
\titleclass{\subparagraph}{straight}
\titleformat{\section}[block]
{\centering\bfseries} % Полужирный, по центру, прописные
{ГЛАВА \thesection.\ } % Нумерация с точкой
{0em} % Расстояние между номером и текстом
{\MakeUppercase} % Код перед заголовком
 % Расстояние после заголовка (2 интервала = 3ex при 1.5)

\titleformat{\subsection}[block]
{\centering\bfseries}
{\thesubsection.\ }
{0em}
{\MakeUppercase}

\titleformat{\subsubsection}[block]
{\centering\bfseries}
{\thesubsubsection.\ }
{0em}
{}

\titleformat{name=\section,numberless}[block]
{\centering\bfseries}
{}
{0em}
{\MakeUppercase}

\titlespacing*{name=\section,numberless}{0pt}{1.3em plus 3pt minus 2pt }{1.3em plus 3pt minus 2pt }


\titleformat{\paragraph}[block]
{\centering}
{\theparagraph. \ }
{0em}
{ }



\titlespacing*{\section}{0pt}{1.3em plus 3pt minus 2pt }{1.3em plus 3pt minus 2pt } % Отступы вокруг заголовка главы
\titlespacing*{\subsection}{0pt}{1.3em plus 3pt minus 2pt }{1.3em plus 3pt minus 2pt } % Отступы вокруг заголовка параграфа
\titlespacing*{\subsubsection}{0pt}{1.3ex plus 3pt minus 2pt }{1.3em plus 3pt minus 2pt } % Отступы вокруг заголовка параграфа
\titlespacing*{\paragraph}{0pt}{1em plus 3pt minus 2pt }{1em plus 3pt minus 2pt } % Отступы вокруг заголовка параграфа

%===================================================
% Стиль страниы и нумерация
%===================================================

\usepackage{fancyhdr}
\pagestyle{fancy}
\fancyfoot{} %пустой нижний колонтитул
\fancyhead{}%Отчищаем верхний колонтитул
\fancyhead[C]{\thepage} % Номер страницы по центру вверху
\renewcommand{\headrulewidth}{0pt} % Убираем черту


%\setlength{\headsep}{10pt}
%\setlength{\voffset}{1cm}
\setlength{\topmargin}{-1in + 13mm}
\setlength{\headheight}{4mm}
\setlength{\headsep}{3mm}
%наскоклько верний колонтиту выше страницы
%\setlength{\headheight}{\heightof{9}}
%\setlength


%переопределяем стиль plain который используется поумолчанию для страниц начала глав
\fancypagestyle{plain}{%
	\fancyfoot{} %пустой нижний колонтитул
	\fancyhead{}%Отчищаем верхний колонтитул
	\fancyhead[C]{\thepage} % Номер страницы по центру вверху
	\renewcommand{\headrulewidth}{0pt} % Убираем черту
%	\setlength{\headsep}{2ex}%наскоклько верний колонтиту выше страницы
%	\setlength{\headheight}{2ex}
}


%===================================================
% Подписи к рисункам, таблицам, листингу
%===================================================

%--- Настройка меток
\usepackage{caption}

%Разделитель номер рисунка
\DeclareCaptionLabelSeparator{dot}{.~}  
%Разделитель номер таблицы         
\DeclareCaptionLabelSeparator{dot_tbl}{.\\}            

%Подпись рис. центр
\captionsetup[figure]{
	name=Рис.,
	justification=centering,
	labelsep=dot, 
	format=plain,
	font={onehalfspacing},
	%skip=0ex, 
	belowskip=-2ex plus 4pt minus 2pt , 
	%aboveskip=1ex plus 4pt minus 2pt 
}                       

%Подпись табл. справа
\captionsetup[table]{
	name=Таблица,
	justification=raggedleft,
	labelsep=dot_tbl, 
	format=plain, 
	singlelinecheck=false,
	font={onehalfspacing},
	%belowskip=1ex plus 4pt minus 2pt, 
	%aboveskip=1ex  plus 4pt minus 2pt 
} 

%Подпись листинг. справа
\captionsetup[listing]{
	justification=raggedleft,
	labelsep=dot_tbl, 
	format=plain, 
	position=above, 
	singlelinecheck=false,
	font={onehalfspacing}, 
	%belowskip=1 ex plus 4pt minus 2pt , 
	aboveskip=-1.4ex  plus 4pt minus 2pt 
} 


%В рисунках и таблиц приложений используем префикс П
\appto\appendix{%
	\renewcommand{\thefigure}{\thesection.\arabic{figure}}
	\renewcommand{\thetable}{\thesection.\arabic{table}}
	\renewcommand{\thelisting}{\thesection.\arabic{listing}}
}


\renewcommand{\thefigure}{\thesection.\arabic{figure}}
\renewcommand{\thetable}{\thesection.\arabic{table}}


%привязываем счетки рисунков и таблиц к секциям (главам)
\usepackage{chngcntr}
\counterwithin{figure}{section}
\counterwithin{table}{section}

% Ссылки на рисунок и таблицу в тексте (требования к оформлению ВКР 2026)
\newcommand{\figref}[1]{Рис.~\ref{#1}}
\newcommand{\tabref}[1]{Таблица~\ref{#1}}


%===================================================
% Список литературы
%===================================================

\usepackage[style=gost-numeric, 
		sorting=none, 
		%defernumbers=true,
		]{biblatex}  
\addbibresource{vrkbibliograpgy.bib}
\urlstyle{same} 



\DeclareDelimFormat{multicitedelim}{\addcomma\space}
%\DeclareDelimFormat{compcitedelim}{\addcomma\space}

\defbibenvironment{bibliography}
{\list
	{\printfield{labelnumber}.\stepcounter{vkrTotalbooks}}
	{\toggletrue{bbx:gostbibliography}% %включаем ГОСТ 2018
		\renewcommand*{\revsdnamepunct}{\addcomma}%
		\renewcommand*{\labelnamepunct}{\addperiod\space}%
		\setlength{\labelwidth}{0pt}%
		\setlength{\leftmargin}{0em}%         % Красная строка для номера
		\setlength{\itemindent}{10.5mm}%        % Компенсируем отступ
		\setlength{\labelsep}{0.4em}%
		\setlength{\listparindent}{0pt}%
		\setlength{\itemsep}{\bibitemsep}%
		\setlength{\parsep}{\bibparsep}%
		\renewcommand*{\makelabel}[1]{\hss##1}}}
{\endlist}
{\item}

\renewcommand*{\mkgostheading}[1]{#1} %убираем курсив в спсике литературы


%===================================================
% Настройка оглавления
%===================================================


\usepackage{tocloft} % Для настройки оглавления

%Заголовок оглавления
\renewcommand{\cfttoctitlefont}{\hfill\bfseries}
\renewcommand{\cftaftertoctitle}{\hfill\vspace{1em}}

\renewcommand{\cftpnumalign}{r}


\cftsetrmarg{3ex}
\cftsetpnumwidth{3ex}

\renewcommand{\cftdotsep}{2}

\renewcommand{\cftsecpresnum}{ГЛАВА }
\setlength{\cftsecnumwidth}{\widthof{ГЛАВА 10.\ }}
\renewcommand{\cftsecaftersnum}{. }
\renewcommand{\cftsecleader}{\cftdotfill{\cftdotsep}}
\renewcommand{\cftsecfont}{\normalfont}        % обычный шрифт
\renewcommand{\cftsecpagefont}{\normalfont}    % шрифт номеров страниц
\setlength{\cftbeforesecskip}{0.0em} 


\setlength{\cftsubsecnumwidth}{\widthof{0.0.\ }}
\renewcommand{\cftsubsecaftersnum}{. } % Точка после номера параграфа
\renewcommand{\cftsubsecleader}{\cftdotfill{\cftdotsep}}
\renewcommand{\cftsubsecfont}{\normalfont}        % обычный шрифт
\renewcommand{\cftsubsecpagefont}{\normalfont}    % шрифт номеров страниц
\setlength{\cftbeforesubsecskip}{0.0em} 

\setlength{\cftsubsubsecnumwidth}{\widthof{0.0.0.\ }}
\renewcommand{\cftsubsubsecaftersnum}{. } % Точка после номера параграфа
\renewcommand{\cftsubsubsecleader}{\cftdotfill{\cftdotsep}}
\renewcommand{\cftsubsubsecfont}{\normalfont}        % обычный шрифт
\renewcommand{\cftsubsubsecpagefont}{\normalfont}    % шрифт номеров страниц
\setlength{\cftbeforesubsubsecskip}{0.0em} 

%Главы капсом
%https://tex.stackexchange.com/questions/156916/how-to-make-section-name-uppercase-in-toc
\makeatletter
%\patchcmd{\l@subsection}{#1}{\MakeUppercase{#1}}{}{}
%\patchcmd{\l@section}{#1}{\MakeUppercase{#1}}{}{}
\makeatother

% Уменьшить отступ перед номерами страниц

%===================================================
% Настройка приложений
%===================================================


% Переопределяем формат для приложений

\appto\appendix{%
	
	%\renewcommand{\thesection}{\arabic{section}} 
	
	\renewcommand{\thesection}{\Asbuk{section}} %почему то опять вернули буквы в приложения
	
	\titleformat{\section}[block]
	{\bfseries\raggedleft} % Полужирный, по центру, прописные
	{Приложение \thesection.} % Нумерация с точкой
	{0em} % Расстояние между номером и текстом
	{\\} % Код перед заголовком
	[\vspace{0ex}] 
	
	\titleformat{\subsection}[block]
	{\centering\bfseries}
	{\thesubsection.\ }
	{0em}
	{}
	
	{ }% Настройки оглавления для приложений
	\addtocontents{toc}{%
		\protect\renewcommand{\protect\cftsecpresnum}{Приложение }%
		\protect\setlength{\protect\cftsecnumwidth}{\widthof{Приложение М.\ }}%
		\protect\renewcommand{\protect\cftsecaftersnum}{.}%
	}
}

%Команды для разделов для прриложений
%Имеют номер, но не светятся в оглавлении
\newcommand{\hiddensubsection}[1]{%
	\addtocounter{subsection}{1}%
	\subsection*{\thesubsection. #1}%	
}

\newcommand{\hiddensubsubsection}[1]{%
	\addtocounter{subsubsection}{1}%
	\subsubsection*{\thesubsubsection. #1}%	
}



%===================================================
% Листинги кода
%===================================================


\usepackage[newfloat]{minted}
\setminted{linenos=false,
	style=vs,
	fontsize=\fontsize{14pt}{14pt},
	tabsize=3,
	frame=lines}
%\usepackage{listings}

%\ifn

\counterwithin{listing}{section}
\renewcommand{\thelisting}{\thesection.\arabic{listing}}
%\renewcommand{\listingscaption}{Листинг}  %для Minted без newfloat

%===================================================
% Настройка станартного положения рисунков
%===================================================

\usepackage{float}

% 1. Задаём [ht] для всех figure по умолчанию
\floatplacement{figure}{hbt} %в месте вызова или внизу станицы
\floatplacement{table}{hbt} %в месте вызова или внизу станицы
\floatplacement{listing}{hbt} %в месте вызова или внизу станицы

\setcounter{bottomnumber}{3}
\setcounter{topnumber}{2}

%счеткики для подсчета колчества рисунков и таблицу

\newcounter{vkrTotalfigures}
\newcounter{vkrTotaltables}

\AtBeginEnvironment{figure}{\stepcounter{vkrTotalfigures}}
\AtBeginEnvironment{table}{\stepcounter{vkrTotaltables}}

%расстояние между float (таблицы, рисунки) и текстом
\setlength{\floatsep}{1.6em plus 2pt minus 2pt}     % между float'ами
\setlength{\textfloatsep}{1.6em plus 2pt minus 2pt} % между float'ами и текстом
\setlength{\intextsep}{1.6em plus 2pt minus 2pt}    % для float'ов внутри текста


\SetupFloatingEnvironment{listing}{name=Листинг} 
%\renewcommand{\lstlistingname}{Листинг}

%для длинных листингов
\newenvironment{longlisting}{\vspace{\intextsep}\captionsetup{type=listing}}{}
\newenvironment{VKRlonglisting}{\vspace{\intextsep}\captionsetup{type=listing}}{}

%для длинных таблиц
\newenvironment{VKRlongtable}{\vspace{\intextsep-1ex}\stepcounter{vkrTotaltables}\begingroup\captionsetup{type=table}\def\insidetable{true}}{\endgroup}

\usepackage{placeins}  % Для команды \FloatBarrier

%===================================================
% Списки
%===================================================


\usepackage{enumitem} %упарвление списками

\setlist[enumerate]{nosep, leftmargin=1cm }  
\setlist[itemize]{nosep, leftmargin=1cm}

%\setlist[enumerate,1]{leftmargin=1.25cm +2ex + \labelsep }  
%\setlist[itemize,1]{leftmargin=1.25cm +2ex + \labelsep }

\setlist[enumerate,1]{leftmargin = \parindent }  
\setlist[itemize,1]{leftmargin = \parindent  }

% русские буквы в счетчике https://tex.stackexchange.com/questions/720496/how-to-make-cyrillic-alphabet-numbering-in-enumerate-with-enumitem-and-polygloss
\makeatletter
\AddEnumerateCounter{\asbuk}{\russian@asbuk@alph}{щ} % section 8.1
\makeatother



%===================================================
% Прочие плагины
%===================================================



\usepackage{lipsum}%для отладки текста, команда \lipsum
\usepackage{clipboard}
\usepackage{ifthen}

\usepackage{marginnote} %заметки на полях
\reversemarginpar 
\setlength{\marginparwidth}{25mm}
\usepackage[textsize=tiny]{todonotes}

\usepackage{refcount}  %всяукие команды для работы с ссылками, например \getpagerefnumber
\usepackage{metalogo}

%===================================================
% Отладка
%===================================================

\ifDebug
	\usepackage{layout} %Позволяет использовать команду \layout для печати структуры страницы
	\usepackage{layouts} %еще мног о команд для откладки
	\usepackage{showframe} %Рамки фокруг листа
	
	\AtBeginDocument{%
		\input{debug_page.tex}%
	}
	
\fi


%===================================================
% Добавляйте сюда ваши собственные пакеты и команды
%===================================================

% Ссылка на формулу в тексте (требования к оформлению ВКР 2026)
\newcommand{\formref}[1]{формула~\eqref{#1}}

% Блок расшифровки символов под формулой: все строки по левому краю
\newcommand{\formwhere}[1]{%
	\par\vspace{0.25\baselineskip}%
	\noindent
	#1%
	\par\vspace{0.25\baselineskip}%
}

\newcommand{\formwherenext}[2]{%
	\par\noindent #1~--- #2%
}



%===================================================
% Пакет hyperref. 
% Должен подключатся последним.
% Он необходим для работы с ссылками
% И меняет работу других пакетов
%===================================================


\usepackage{hyperref}
\hypersetup{hidelinks} %отключить посдветку ссылок
%\usepackage{bookmark} %для заклкдок



%===================================================
% Кастомные команды
%===================================================

%===========================
% Запись/чтение из aux

\makeatletter
% Команда для записи в .aux
\newcommand{\writetoaux}[2]{%
	
	\immediate\write\@auxout{%
		\string\global\string\@namedef{vkr#1}{#2}%
	}%
	
}
\makeatother

% Команда для чтения из .aux
\newcommand{\readfromaux}[1]{%
	\ifcsname vkr#1\endcsname
	\csname vkr#1\endcsname
	\else
	*Счетчик не определен или требуется перекомпиляция*%
	\fi
}

%===========================

%Красная маленькая заметка на полях

\newcommand{\vkrnote}[2][]{%
	\marginnote[{\fontsize{9pt}{9pt}\selectfont\fussy\color{red}#1}]%
	{{\fontsize{9pt}{9pt}\selectfont\setlength{\baselineskip}{9pt}\fussy\color{red}#2}}%
}


%===========================


%разбивка  строки на слова-строки
\ExplSyntaxOn
\NewDocumentCommand{\wordsperline}{m}
{
	\seq_set_split:Nnn \l_tmpa_seq { ~ } { #1 }
	\seq_use:Nn \l_tmpa_seq { \\ }
}
\ExplSyntaxOff

%Станкиновская шапка

\newcommand\vkrStandartHead%
{
	\begingroup
	\fontsize{12pt}{10pt}\selectfont
	{
		\centering
		\bfseries\noindent
		\includegraphics[width=54mm]{stankin-logo.pdf}\\
		МИНОБРНАУКИ РОССИИ\\
		федеральное государственное автономное образовательное учреждение\\
		высшего образования\\
		«Московский государственный технологический университет «СТАНКИН»\\
		(ФГАОУ ВО <<МГТУ <<СТАНКИН>>) \\
	}
	
	
	
	\vspace{-0.75em}\noindent\rule{\textwidth}{0.4pt}\vspace{2ex}
	
	\noindent	
%	\begin{minipage}[t]{0.3\textwidth}
%		\bfseries
%		Институт\\
%		Информационных\\
%		Технологий\\	
%	\end{minipage}
%	\hfill
%	\begin{minipage}[t]{\widthof{\bfseries Вычислительных систем}+1pt}
%		\bfseries
%		Кафедра\\
%		Информационных\\
%		Технологий и\\ 
%		Вычислительных систем\\	
%	\end{minipage}
	\begin{minipage}[t]{0.4\textwidth}
		\bfseries
		Институт\\\expandafter\wordsperline\expandafter{\vkrInstitute}
	\end{minipage}
	\hfill
	\begin{varwidth}[t]{0.3\linewidth}
		\bfseries\raggedright
		Кафедра\\
		\vkrKafedra
	\end{varwidth}
	\endgroup
}

\def\emptystr{}
%===========================
% Макросы для вставки должности регалий о руководителе или консультанте


%должность, степ. , зван.
\def\insertPersonInfo#1{%
	\edef\tmpZ{\csname vkr#1Position\endcsname}%
	\edef\tmpA{\csname vkr#1DegreeShort\endcsname}%
	\edef\tmpB{\csname vkr#1TitleShort\endcsname}%
	\tmpZ%
	\ifx\tmpA\empty\else%
	,\space\tmpA%
	\fi%
	\ifx\tmpB\empty\else%
	,\space\tmpB%
	\fi%	
}

%степ. , зван., если нет то должность
\def\insertPersonInfoA#1{%
	
	\edef\tmpZ{\csname vkr#1Position\endcsname}%
	\edef\tmpA{\csname vkr#1DegreeShort\endcsname}%
	\edef\tmpB{\csname vkr#1TitleShort\endcsname}%
	\ifx\tmpA\emptystr%
	\ifx\tmpB\empty%	
	\tmpZ%
	\else%
	Error In #1 info!%
	\fi%		
	\else%
	\tmpA%
	\ifx\tmpB\empty\else%
	,\space\tmpB%
	\fi%
	\fi%
}

%должность каф. , степ. , зван.
\def\insertPersonInfoB#1{%
	\edef\tmpZ{\csname vkr#1Position\endcsname}%
	\edef\tmpA{\csname vkr#1DegreeShort\endcsname}%
	\edef\tmpB{\csname vkr#1TitleShort\endcsname}%
	\tmpZ\space каф. \vkrKafedraShort%
	\ifx\tmpA\empty\else%
	,\space\tmpA%
	\fi%
	\ifx\tmpB\empty\else%
	,\space\tmpB%
	\fi%	
}


%===========================
%===========================
%===========================
%===========================
%===========================
%===========================
%===========================
%===========================
%===========================
%===========================
%===========================
%===========================
%===========================